\subsection{Ising twist masses at multiplicity eight}
\label{subsec:round4-m8-ising}
The two multiplicity-eight rings $R_{6,8,311}$ and $R_{6,8,1660}$ have
formal codegrees with characteristic polynomial
\[
 (t-4)^2(t^4-576t^3+11520t^2-73728t+147456).
\]
They contain Ising basis subcategories, with labels $(g,\sigma)=(b_4,b_3)$
and $(b_1,b_4)$, respectively.  Their different duality involutions
remain part of the full tensor input; the common polynomial does not
identify the rings.

Here the entire character field embeds in $\QQ(\zeta_{24})$.  A
pseudo-unitary Ising category has twists $1,-1,\xi$, where $\xi$ is a
primitive sixteenth root.  Its center has four $\sigma$-restriction
objects, with twists $\xi,-\xi,\xi^{-1},-\xi^{-1}$.  On taking normalized
trace to $\QQ(\zeta_{24})$, their height-three vectors become
\[
 (0,\eta,0),\qquad (0,\eta^{-1},0),\qquad \eta=\xi^2,
\]
with each vector occurring for two of the four objects.  The sign
$\epsilon=\nu_2(\sigma)$ determines the unordered pair:
\[
 \eta+\eta^{-1}=\epsilon\sqrt2.
\]
Both possible signs are retained whenever permitted by the second
indicator conditions.

Corollaries~\ref{cor:fibonacci-ising-projection} and
\ref{cor:subcategory-twist-masses} now give the following simultaneous
constraints.  Every low column satisfies $c_\sigma=0$ and
$c_g\leq c_0$; moreover
\[
 \sum_{\text{low }Z}d_Zc_g(Z)=D/2.
\]
A high column with $c_g>0$ has twist $-1$ and $c_\sigma=0$.
A high column with $c_\sigma>0$ has $c_g=0$ and one of the two displayed
projected vectors.  For each of those vectors the exact dimension mass
is $D\sqrt2/2$.  These equations are imposed together with the complete
induction Gram, full-character reciprocity, Galois balance, and cyclic
rotation equations.  They do not assume that the Ising braiding
extends to the ambient category.

\begin{theorem}
\label{thm:round4-311}
The fusion ring $R_{6,8,311}$ has no categorification over $\CC$.
\end{theorem}
\begin{proof}
The unrestricted signed reduction has $96$ quotients and eight sign
orbits.  One orbit is gauge equivalent to the original ring, and $13$
exact dimension-character checks reject all other possible categorical
dimension components.  The surviving component is the
Frobenius--Perron Galois orbit, so it suffices to exclude a
pseudo-unitary categorification.

The complete low induction domain has two frames.  Each has two
character labellings and four allowed second-indicator tuples.  With
$F=24$ and height three, the complete order bound is $144$ and there are
$53$ normalized trace vectors.  The first frame gives systems of $387$
equations in $802$ variables; the second gives $1017$ equations in
$1896$ variables.  Each of the $16$ branches has an exact rational
separating vector.  The direct certificate is
\texttt{research/round4/m8/R6\_8\_311\_certificate.json}; the
full-tensor verifier is
\texttt{software/round4\_m8\_ising.py}.
\end{proof}

The same criterion is implemented uniformly and screened on all ten
frontier inputs and all $54$ realized controls.  Three frontier inputs
contain an Ising basis subcategory: the two rings just discussed and
$R_{6,4,91}$, treated in its own campaign.  Six realized rings contain
one.  Five are outside the present full-character-field implementation,
and the applicable control $R_{6,2,94}$ survives the necessary system.

For $R_{6,8,1660}$, the unrestricted signed reduction has $640$
quotients in $24$ sign orbits and $58$ nongauge dimension-character
rejections. It proves sphericality and a pseudo-unitary Galois reduction.
Morita reconstruction over $1+g$ returns the same based ring, so it
supplies no transfer to the excluded non-self-dual partner. The complete
current profile enumeration and actual-root twist constraints are given
in the frontier analysis below, replacing the earlier bounded searches.
